\section{Discussion}\label{sec:discussion}

\subsection{Implications for field operation}\label{sec:field}

The scheme that did best is also the simplest to run. It needs the
silica grab samples that field chemistry programs already take, here
every 8~h, one scaling-rate measurement per day, and a daily calculation
that takes well under a second on a laptop. The correction is one
number per day: the measured rate minus the model rate at the measured
silica concentration and the current temperature and dose, filtered
with a gain of one half. Two choices need care. The first is the
anchor. An operator who compares the measured rate with the design-case
prediction, a natural reference when the design case is the document of
record, will read the design margin as model error. The second is the
back-off. Under sixfold faster kinetics a gain of $\kappa = 1$ on the
correction's own prediction-error standard deviation exchanged
\CT{boCA}{u}{abs} points of power gain for \CT{boCA}{v}{abs} fewer
violation days per run, while without systematic kinetic error it cost
\NUM{boCAB} points and bought nothing. A site without kinetic data of
its own should therefore start with $\kappa$ near one and lower it as
the history of the correction shows the model error to be small.
Because the safe response to a kinetic error is a few kelvin of
reinjection temperature (Section~\ref{sec:operating}), the set-point
changes the scheme asks for are of a size an operator would treat as
routine.

A field trial can reuse the pilot designs of Philippine inhibitor
programs, in which a side stream of treated brine runs through an
insulated test line with inspection spools and retention vessels while
temperature, pressure and silica are monitored~\citep{see2010wgc}. Run
at the optimized set-points, such a line would test the policy and,
through its spools, calibrate the scaling limit against wall growth.

\subsection{Scaling limit and wall growth}\label{sec:units}

The constraint is a polymerization (monomer-depletion) rate, not a rate
of wall growth. The molecular-deposition treatment of
\citet{weres1981kinetics}, transcribed under the same integrity protocol
(its printed checkpoints reproduce, including the 3.0~kcal/mol effective
activation energy; one internal inconsistency in its Eq.~[4c] is
documented and not silently repaired), converts
$r_{\max} = 0.5$~g\,t$^{-1}$\,min$^{-1}$ into a wall growth of about
0.3~mm/yr at the optimized operating point and 0.04--1.6~mm/yr across
90--175~$^{\circ}$C. That is a lower bound on scale growth. Field
deposits also grow by the aggregation and cementation of colloidal
silica~\citep{vandenheuvel2018understanding}, and inhibitor-treated
Philippine lines have shown 2.4--11.9~mm/yr in a 0.5~ppm test at Leyte
and up to about 20~mm/yr projected along a hot-injection pilot line at
Bacman~\citep{baltazar2014grc,see2010wgc}. The limit used here is
therefore a placeholder, to be calibrated against spool or coupon data
at the site. In the sensitivity study a fourfold band of $r_{\max}$
moves the optimal temperature far less than the kinetic uncertainty
does (total-order index \NUM{SobSTTrmax} against \NUM{SobSTT}).

\subsection{Limitations}\label{sec:limits}

The closed-loop results are in silico. The plant is simulated, although
every physical ingredient, from solubility and kinetics to the brine
envelope and the inhibitor response, comes from primary literature and,
wherever independent data exist, was tested against them
(Section~\ref{sec:labtest}). Component-wise validation on real data
combined with a closed loop under parameterized mismatch is an
intermediate step between simulation-only studies and a field trial,
not a substitute for the trial.

The ranking among the measured-anchored corrections holds for a line
held near one operating point, with one label per day, weekly refits of
the learned models, and the particular regressors and uncertainty
estimates used here. A line whose brine steps when wells are switched
would favour a correction that depends on the state over one that is
merely current, and that case was not tested. The Gaussian-process arm
compares a regressor, not Gaussian-process modifier adaptation, whose
trust regions and acquisition functions might change its curve.

Violations are counted on the realized plant. Every feedback policy, by
tightening or relaxing the constraint from what the plant returns,
trades the worst-case guarantee of Eq.~\eqref{eq:rto} for performance
at the realized brine. Under the sixfold stress a policy that kept the
full envelope guarantee would have to run hotter than fixed practice,
which is what the log-space correction at the measured anchor does; the
additive corrections are safe on the realized plant but would
under-correct if the brine moved to the worst-case vertex.

The chemistry is near neutral. High-pH brines need silica ionization
corrections that are deferred here (the speciation layer exists in the
repository but is not used in the loop), and acid dosing, the third
lever of Philippine practice~\citep{remoroza2008ph}, is not modeled.
Economic figures are scenario values with declared ranges.

Finally, the anchoring ablation exists because an initial analysis used
the design-case anchor for its scalar baseline and concluded, wrongly,
that the hybrid loop was the only safe policy. The ablation, the
measured-anchored scalar family, the Gaussian-process comparator, the
gain sweep and fourteen of the twenty plants were added after that
error was found. The effect of the anchor, an order of magnitude in
violation days with intervals far from zero and a mechanism reproduced
by a separate diagnostic, is secure; the finer comparisons among the
corrections were not registered in advance and would merit replication.

\subsection{Beyond silica}\label{sec:general}

The anchoring result does not depend on silica chemistry. Whenever an
operating point is optimized against a worst-case envelope of
composition and corrected from plant feedback, the feedback has to be
compared with the model at the realized composition, because at the
design case it carries the robustness margin as well as the model
error. The steeper the scaling law, the larger the margin in rate terms
and the worse the confusion: here the design-case brine, at 911~mg/kg,
is only 7\% more concentrated than the nominal operating brine, yet the
model rates at the two differ by a median factor of \NUM{AnMarginmed}. The same care applies to calcite, stibnite or barite limits set
from envelope calculations~\citep{tonkul2021characterization,%
omrani2025probabilistic}. The second lesson is more modest: on this
problem a better estimator did not make a better controller, and the
simplest correction that used the right measurement was as good as any.
